\chapter*{Preface}
\addcontentsline{toc}{chapter}{Preface}

\hayashi{} was born out of a concrete frustration: existing tools for applied
econometrics force the researcher to choose between expressiveness and
performance, between a syntax close to statistical notation and the ability
to automate complex analyses.

Stata has the right syntax, but it is proprietary, closed, and not a
general-purpose language. R has the right ecosystem, but its syntax is
inconsistent and its data model is surprising for those coming from Stata.
Python has the right flexibility, but its statistical packages treat
econometrics as a special case of machine learning.

\hayashi{} is an interpreted, general-purpose language built specifically
for econometric analysis. The syntax pays homage to Stata where it matters
and adopts modern idioms where Stata shows its age. The engine is written
in Rust — no virtual machine, no garbage collector, no perceptible startup
time.

\bigskip

This book documents both the language and the econometric core: from the
type structure and control flow to time series estimators, causal inference,
and advanced methods. It is a snapshot of the current state of the project
— a living document that evolves with the software.

\bigskip
\noindent\rule{\textwidth}{0.4pt}
\bigskip

\section*{A project by one developer — and by the entire community}

\hayashi{} is, to date, the work of a single developer. This has practical
consequences the reader should know about.

On the positive side, there is design coherence: every syntax decision,
every estimator interface, every error message was made with the same set
of values in mind. There are no committees, no bikeshedding, no
compatibility layers accumulated over decades. The code is lean because
nobody added anything ``just to keep it''.

On the limitation side, one pair of eyes — however careful — does not see
everything. The project relies on three layers of verification: unit tests
covering the language primitives; integration scripts that exercise
\textsc{Hayashi} as a black box; and DGP (data-generating process) scripts
for each estimator, where the true parameter is known and recovery can be
verified numerically. Even so, edge cases exist, and some have certainly
escaped detection.

\bigskip

\noindent\textbf{License.} \hayashi{} is free software, distributed under
the terms of the \textbf{GNU General Public License version~3} (GPL~v3).
This means you may use, study, modify, and redistribute the program —
including for commercial purposes — provided that any derivative work is
equally free. The source code is public and belongs, in the most literal
sense, to whoever uses it.

The choice of GPL~v3 is not accidental. Proprietary econometric tools
create institutional dependency: universities and research labs become
hostage to licenses, pricing policies, and corporate decisions about what
the software can or cannot do. Free software breaks that cycle.
\hayashi{} should be accessible to any researcher anywhere in the world,
regardless of funding.

\bigskip

\noindent\textbf{Contributors and testers are welcome.} If you found an
unexpected result, an estimator that fails to converge when it should, an
obscure error message, or simply a syntax that could be clearer — that is
valuable information. Open an issue in the repository. If you want to go
further, a well-tested fix or a new documentation chapter is as legitimate
a contribution as any line of Rust code.

The project especially needs people willing to test \hayashi{} on real
data: time series with seasonality, unbalanced panels, models with
multicollinearity, data with missing values. Those are the cases that DGP
scripts do not cover, because synthetic data is, by construction,
well-behaved.

\hayashi{} only becomes robust through use. And honest use — which includes
reporting what breaks — is the most direct way to contribute.

\bigskip

\noindent\textit{June 2026.}
